%%%%%%%%%%%%%%%%%%%%%%%%%%%%%%%%%%%%%%%%%%%%%%%%%%%%%%%%%
%% © DRL CSAIL MIT 2022 - All RIGHTS RESERVED.
%% AUTHORS: RAMIN HASANI and ALEXANDER AMINI
%%
%% License: Creative Commons license
%% Attribution-ShareAlike 4.0 International (CC BY-SA 4.0)
%%%%%%%%%%%%%%%%%%%%%%%%%%%%%%%%%%%%%%%%%%%%%%%%%%%%%%%%%

\documentclass{MITcsail}
\usepackage[numbers]{natbib}
\usepackage{amsthm}
\usepackage{graphicx}
\usepackage{booktabs}
\usepackage{textcomp}   % \textperthousand
\usepackage{amsmath}

\newcommand{\permil}{\textperthousand}

\title{COS 568 Final Project: Task 2, Milestone 2}

\author{Anderson Lee}

\begin{document}

\maketitle
\thispagestyle{firstpagestyle}

%------------------------------------------------------------------
\section{Experimental Setup}
%------------------------------------------------------------------
Experiments were run on the Della cluster with 12 CPUs per task, 1 task per
node, and 64\,GB of memory per task.  All benchmarks used the Facebook dataset
(\texttt{fb\_100M\_public\_uint64}) with two mixed workloads: a
\emph{lookup-heavy} mix (10\% inserts, 90\% lookups) and an
\emph{insert-heavy} mix (90\% inserts, 10\% lookups).  Each configuration was
repeated $r{=}3$ times; throughput figures are averages of those three runs.

%------------------------------------------------------------------
\section{Hybrid Design}
%------------------------------------------------------------------
The hybrid index combines \textbf{LIPP} as a read-optimised main store with
\textbf{Dynamic PGM (DPGM)} as a write-absorbing insertion buffer.

\paragraph{Build phase.}
The initial (bulk-loaded) dataset is loaded entirely into LIPP via
\texttt{bulk\_load}.  DPGM starts empty.

\paragraph{Insert.}
Every new key--value pair is inserted into DPGM.  After each insert, the
hybrid checks whether DPGM has accumulated
$\lfloor N_0 \times t / 1000 \rfloor$ keys, where $N_0$ is the initial key
count and $t$ is the \texttt{flush\_threshold\_per\_mille} template
parameter.  When the threshold is reached, a synchronous \emph{flush}
drains all DPGM entries individually into LIPP and resets DPGM to an empty
state.

\paragraph{Lookup.}
Equality lookups probe DPGM first (the buffer is typically small), then fall
back to LIPP on a miss.

\paragraph{Why per-mille thresholds.}
With $N_0 = 100\text{M}$ initial keys and a 2\,M-operation workload
containing at most 1.8\,M inserts, using integer \emph{percent} thresholds
would set the minimum flush trigger at 1\,M keys---meaning the lookup-heavy
workload (only $\sim$200\,k inserts) would \emph{never} flush.  Per-mille
(‰) thresholds keep the same template-parameter name visible in CSV
variants while allowing realistic triggers (e.g.\ $2\permil \times 100\text{M}
= 200\text{k}$ keys).

%------------------------------------------------------------------
\section{Implementation Notes}
%------------------------------------------------------------------
\paragraph{Naive synchronous flush.}
Rather than migrating data asynchronously (Milestone~3), this implementation
performs a blocking flush: the \texttt{Flush()} method iterates over all
DPGM entries via \texttt{for\_each} and inserts each key individually into
LIPP.  This is correct and simple, but it blocks all operations for the
duration of the flush.

\paragraph{LIPP does not support bulk re-insertion.}
After its initial \texttt{bulk\_load}, LIPP accepts new keys only via its
single-key \texttt{insert} interface.  Flushing therefore must iterate one
key at a time; there is no faster batch path available.

\paragraph{Ordered vs.\ unordered traversal.}
The flush uses DPGM's \texttt{for\_each} (level-by-level, unordered)
rather than the sorted iterator, because insertion order into LIPP does not
affect correctness and the sorted path exposed a bug in earlier iterator
versions.

%------------------------------------------------------------------
\section{Hyperparameter Selection}
%------------------------------------------------------------------

\paragraph{DPGM error $\varepsilon$.}
On the Facebook mixed workloads, three values of the PGM error parameter
$\varepsilon \in \{64, 128, 256\}$ were evaluated.  For both workloads,
$\varepsilon{=}64$ achieved the highest average throughput
(Table~\ref{tab:results}):
$1.089$\,Mops/s on the lookup-heavy workload and $3.590$\,Mops/s on the
insert-heavy workload.  Smaller $\varepsilon$ tightens the prediction interval,
reducing the local search window and improving lookup speed; the insert cost
increase is marginal.

\paragraph{Hybrid flush threshold.}
Three threshold candidates were evaluated per workload (see
Table~\ref{tab:results}).

\begin{itemize}
  \item \emph{Lookup-heavy (10\% inserts, $\sim$200\,k total inserts).}
    A threshold of $1\permil$ (100\,k keys) triggers roughly two flushes
    during the workload.  This keeps the DPGM buffer small almost always,
    so most lookups land in LIPP directly with full learned-model speed.
    The best average throughput was $1.246$\,Mops/s at
    $(\varepsilon{=}128,\; t{=}1\permil)$.

  \item \emph{Insert-heavy (90\% inserts, $\sim$1.8\,M total inserts).}
    A threshold of $10\permil$ (1\,M keys) triggers approximately one to
    two flushes.  Infrequent flushing keeps the amortised flush cost
    per-insert low, while the buffer remains manageable in size.
    The best average throughput was $2.352$\,Mops/s at
    $(\varepsilon{=}128,\; t{=}10\permil)$.
\end{itemize}

Notably, the best hybrid $\varepsilon$ is 128 rather than 64.  When DPGM
acts as a small buffer (and is frequently reset), the marginal benefit of a
tighter search window is smaller, while a coarser $\varepsilon$ reduces the
number of model segments and speeds up DPGM inserts.

%------------------------------------------------------------------
\section{Results}
%------------------------------------------------------------------

\begin{table}[h]
\centering
\caption{Benchmark results on the Facebook dataset (\texttt{fb\_100M}), 3 runs each.
  Throughput is the average mixed Mops/s $\pm$ sample standard deviation.
  Index size is measured at the end of the run.
  Best hyperparameters selected by highest average throughput per workload.
  Hybrid uses $\varepsilon{=}128$, flush $1\permil$ for the lookup-heavy workload
  and $\varepsilon{=}128$, flush $10\permil$ for the insert-heavy workload.}
\label{tab:results}
\smallskip
\begin{tabular}{llrr}
\toprule
\textbf{Index} & \textbf{Workload} & \textbf{Throughput (Mops/s)} & \textbf{Index Size (MB)} \\
\midrule
DPGM ($\varepsilon{=}64$) & 90\% Lookup, 10\% Insert & $1.089 \pm 0.006$ & 1626 \\
LIPP & 90\% Lookup, 10\% Insert & \textbf{$12.816 \pm 4.086$} & 12113 \\
Hybrid ($\varepsilon{=}128$, $1\permil$) & 90\% Lookup, 10\% Insert & $1.246 \pm 0.043$ & 12113 \\
\midrule
DPGM ($\varepsilon{=}64$) & 10\% Lookup, 90\% Insert & \textbf{$3.590 \pm 0.021$} & 1626 \\
LIPP & 10\% Lookup, 90\% Insert & $2.199 \pm 0.029$ & 12070 \\
Hybrid ($\varepsilon{=}128$, $10\permil$) & 10\% Lookup, 90\% Insert & $2.352 \pm 0.021$ & 12020 \\
\bottomrule
\end{tabular}
\end{table}

\begin{figure}[h]
  \centering
  \includegraphics[width=\linewidth]{analysis_results/hybrid_10pct_insert_90pct_lookup.png}
  \caption{Lookup-heavy workload (10\% insert, 90\% lookup) on Facebook.
           Left: mixed throughput (Mops/s). Right: index size (MB).}
  \label{fig:lookup-heavy}
\end{figure}

\begin{figure}[h]
  \centering
  \includegraphics[width=\linewidth]{analysis_results/hybrid_90pct_insert_10pct_lookup.png}
  \caption{Insert-heavy workload (90\% insert, 10\% lookup) on Facebook.
           Left: mixed throughput (Mops/s). Right: index size (MB).}
  \label{fig:insert-heavy}
\end{figure}

%------------------------------------------------------------------
\section{Analysis and Discussion}
%------------------------------------------------------------------

\paragraph{Lookup-heavy workload.}
LIPP records the highest \emph{average} throughput at $12.816$\,Mops/s, but
the result comes with an extraordinary standard deviation of
$\pm 4.086$\,Mops/s (coefficient of variation $\approx 32\%$).  The three
individual runs measured $16.54$, $8.44$, and $13.47$\,Mops/s---a nearly
$2\times$ spread.  Because the LIPP structure occupies $\sim$12\,GB, it far
exceeds the processor's last-level cache; throughput is therefore highly
sensitive to the state of the OS page cache and memory-bandwidth competition
from co-scheduled jobs on the shared Della cluster.  DPGM ($1.089$\,Mops/s)
and the Hybrid ($1.246$\,Mops/s) are both compact enough that their working
sets fit comfortably in memory, yielding stable and consistent measurements.
The Hybrid outperforms standalone DPGM by $\approx 14\%$: although 90\% of
operations are lookups, most lookups reach the fast LIPP model directly
because the DPGM buffer is flushed frequently and kept very small.

\paragraph{Insert-heavy workload.}
DPGM leads at $3.590$\,Mops/s because its buffered update strategy absorbs
inserts cheaply.  LIPP drops to $2.199$\,Mops/s---each new key may trigger
conflict resolution and structural rearrangement.  The Hybrid achieves
$2.352$\,Mops/s, beating LIPP by $\approx 7\%$ while falling well short of
pure DPGM.  DPGM's buffer absorbs inserts efficiently, but frequent flushes
(each requiring individual LIPP inserts) impose synchronous pauses that erode
throughput relative to unbuffered DPGM.  Increasing the flush threshold from
$5\permil$ to $10\permil$ helps (fewer but larger flush events), but the
blocking nature of each flush remains the primary bottleneck.

\paragraph{Index size.}
DPGM is by far the most compact structure at $\sim$1\,626\,MB.  LIPP and
the Hybrid are both $\sim$12\,000--12\,100\,MB.  The Hybrid's size is
nearly identical to LIPP's because LIPP stores all $100\text{M}$ initial
keys plus the flushed inserts, and the residual DPGM buffer (reset after
each flush) is negligible.  There is no size advantage to using the Hybrid
over LIPP alone; the Hybrid's value proposition is purely in write
throughput on insert-heavy workloads.

\paragraph{Bottleneck of the naive approach.}
The core limitation of Milestone~2's design is that \texttt{Flush()} is
synchronous: while flushing, no inserts or lookups can proceed.  For
1\,M-key flushes on an insert-heavy workload, each flush is a large
blocking pause.  This motivates an asynchronous, double-buffered design
in Milestone~3, where a background thread migrates DPGM entries to LIPP
while the foreground continues accepting requests into a second DPGM buffer.

%------------------------------------------------------------------
\section{Conclusion}
%------------------------------------------------------------------
The naive Hybrid DPGM + LIPP index behaves as expected on the insert-heavy
workload: it outperforms standalone LIPP by $\approx 7\%$ ($2.352$ vs.
$2.199$\,Mops/s) while remaining below the DPGM baseline.  On the
lookup-heavy workload, LIPP's \emph{average} throughput dominates at
$12.816$\,Mops/s, but the measurement is unreliable (std $= 4.086$\,Mops/s)
due to the structure's large 12\,GB memory footprint interacting with OS page
cache state.  The Hybrid ($1.246$\,Mops/s) is $\approx 14\%$ faster than
standalone DPGM ($1.089$\,Mops/s) on the same workload.  The synchronous flush
is the primary performance ceiling for the insert-heavy case; removing it via
an asynchronous double-buffer scheme is the natural next step.

\end{document}
